\section{Version B: implementation with an exogenous material intensity of capital}\label{sec:mkt:B}

This section sets the instruments of Table~\ref{tab:mkt:instruments} under Version~B, $\Phi^I$
exogenous, and shows that the competitive equilibrium of Section~\ref{sec:mkt:setup:problems}
reproduces the allocation of Problem~\ref{prob:sp:B}. It presupposes Section~\ref{sec:mkt:setup} and
Section~\ref{sec:sp:B}, and nothing from Section~\ref{sec:mkt:A} until
Section~\ref{sec:mkt:B:compare}, which compares the two.


\subsection{What has to be corrected}\label{sec:mkt:B:what}

Under Version~B the accounting is recursive: by~\eqref{eq:sp:B:Omega} the intensity $\Omega$ and the
material intensity of the final good $\Phi^Y$ are computed from the allocation after the fact, and by
Proposition~\ref{prop:sp:B:irrelevance} neither, and no relative intensity $\omega^j$, appears in the
planner's problem. One restriction survives, the mass constraint $\Phi^IG\le R$
of~\eqref{eq:sp:B:mass}, with multiplier $\eta\ge0$ that is zero away from the corner.

The externalities to be corrected are therefore the two of
Section~\ref{sec:mkt:setup:laissezfaire}, and only those.

\smalltitle{(i) Mass drawn into production.} A tonne entering production leaves as waste exactly once,
and collecting, treating and releasing it is costly. The final-good firm pays $P^R$ for the tonne and
nothing for what becomes of it; the mine pays nothing for the overburden its extraction displaces.

\smalltitle{(ii) Emissions.} The collection sector releases $\left(1-\varpi+d^W\varpi\right)W$ tonnes
and recyclers withdraw $d^WR^R$ of that, and neither is priced.

There is no third margin. The composition of final-good spending does not affect how much of the
incoming mass is stored in durables, because $\Phi^I$ is fixed, so there is nothing for an instrument
on output or on consumption to correct. This is the market statement of
Proposition~\ref{prop:sp:B:nowedge}, and Section~\ref{sec:mkt:B:nogate} makes it sharp.


\subsection{The instruments}\label{sec:mkt:B:instruments}

\begin{proposition}[Implementation, Version~B]\label{prop:mkt:B:implementation}
Let $\left\{C,I,N,D,\varpi,K^R,W\right\}$ solve Problem~\ref{prob:sp:B} with associated shadow prices
$\left\{\lambda,q,p^S,p^X,p^P,p^M,p^W,\eta\right\}$, and set
\begin{align}
\tau^R &= p^W-\eta,
&
\tau^{N,S} &= \Omega^{N,S}p^W,
&
\tau^{D,S} &= \Omega^{D,S}p^W,
\label{eq:mkt:B:mass}
\\
\tau^E &= p^P,
&
s^R &= d^Wp^P,
&
\tau^j &= 0,\quad j\in\left\{Y,C,N,D,W\right\},
\label{eq:mkt:B:emission}
\\
\tau^K &= -\Phi^I\left(p^W-p^M-\eta\right),
&&&&
\label{eq:mkt:B:tauK}
\end{align}
with $s^W$ from the zero-profit condition~\eqref{eq:mkt:setup:zeroprofit} and $T$
from~\eqref{eq:mkt:setup:govbudget}. Then the competitive equilibrium of
Definition~\ref{def:mkt:equilibrium} reproduces that allocation, and the equilibrium prices are
\begin{align}
P^R &= p^{N,\mathrm B},
&
B^{\mathrm B} &= p^{N,\mathrm B}+d^Wp^P,
&
P^T &= \alpha B^{\mathrm B},
&
P^K &= F_K,
\notag\\
Q &= q,
&
P^S &= p^S,
&
P^X &= p^X,
&
s^W &= p^W,
\label{eq:mkt:B:prices}
\end{align}
with $p^{N,\mathrm B}$ as in~\eqref{eq:sp:B:Bvalue}, and the market interest rate equals the planner's,
$r=\rho-\dot\lambda/\lambda$ with $\lambda=u'(C)$.
\end{proposition}

\begin{proof}
Term by term against Section~\ref{sec:sp:B:foc}; each step is a usable consistency check on a
numerical solution.

\emph{(1) Materials.} With $\tau^Y=0$ and $\tau^R=p^W-\eta$, the final-good firm's
condition~\eqref{eq:mkt:setup:foc:goodsfirm} gives $P^R=F_R-p^W+\eta$, which by the production
margin~\eqref{eq:sp:B:FR} is $P^R=p^{N,\mathrm B}$. Hence $B^{\mathrm B}=P^R+s^R=p^{N,\mathrm B}+d^Wp^P$, which is the materials
arbitrage~\eqref{eq:sp:B:arbitrage}: the price a recycler receives per tonne recovered equals the
planner's $B^{\mathrm B}$ of~\eqref{eq:sp:B:Bvalue}.

\emph{(2) Recycling.} The recycler's conditions~\eqref{eq:mkt:setup:foc:recycler} then read
$P^T=\alpha B^{\mathrm B}$ and $P^K=a'\!\left(x\right)B^{\mathrm B}$, the latter coinciding with the final-good firm's rental
condition $P^K=F_K$ exactly when the planner's~\eqref{eq:sp:B:foc:KR} holds. The capital market clears
at the planner's $x$ and $K^R$.

\emph{(3) Treatment.} With $P^T=\alpha B^{\mathrm B}$, $\tau^E=p^P$ and $\tau^W=0$, the collector's
condition~\eqref{eq:mkt:setup:foc:varpi} is $\alpha B^{\mathrm B}+\left(1-d^W\right)p^P=dc^T/d\varpi$, which
is~\eqref{eq:sp:B:foc:varpi}; and the zero-profit condition~\eqref{eq:mkt:setup:zeroprofit} becomes
$s^W=c^W\!\left(\varpi\right)+\left(1-\varpi+d^W\varpi\right)p^P-\varpi\alpha B^{\mathrm B}$, which is the
planner's waste condition~\eqref{eq:sp:B:foc:W}: $s^W=p^W$.

\emph{(4) Extraction and exploration.} With $\tau^N=0$ and $\tau^{N,S}=\Omega^{N,S}p^W$, the mine's
condition~\eqref{eq:mkt:setup:foc:mine} reads $P^R=C^N_N+P^S+\Omega^{N,S}p^W$; comparing with
$P^R=p^{N,\mathrm B}$ from step~(1) and the definition of $p^{N,\mathrm B}$ in~\eqref{eq:sp:B:Bvalue} gives $P^S=p^S$. The
exploration condition then gives $P^X=p^X$ against~\eqref{eq:sp:B:foc:D}, and the costate
equations~\eqref{eq:mkt:setup:costate:mine} coincide with~\eqref{eq:sp:B:costate:S}
and~\eqref{eq:sp:B:costate:X} once the discount rates are shown to agree, which is step~(6).

\emph{(5) Consumption and capital.} With $\tau^C=0$ the household's
condition~\eqref{eq:mkt:setup:household-euler} is $u'(C)=\lambda$, which is~\eqref{eq:sp:B:foc:C} and
identifies the household's $\lambda$ with the planner's. With $\tau^K$ as in~\eqref{eq:mkt:B:tauK},
$Q=1+\tau^K=1+\Phi^I\left(p^M-p^W+\eta\right)=q$ by~\eqref{eq:sp:B:foc:I}.

\emph{(6) The interest rate.} Substituting $P^K=F_K$ and $Q=q$
into~\eqref{eq:mkt:setup:household-euler} gives $r=\left(F_K+\dot q\right)/q-\delta$, which is the
planner's capital costate~\eqref{eq:sp:B:costate:K}; and since both rates are defined as
$\rho-\dot\lambda/\lambda$ with the $\lambda$'s identified in step~(5), they coincide. This completes
step~(4), and the household's Euler equation is~\eqref{eq:sp:B:euler}.

\emph{(7) Pollution and the remaining relations.} $P$ is taken as given by household and firms, so
$p^P$ has no private counterpart and enters only through $\tau^E$; its law of
motion~\eqref{eq:sp:B:costate:P} is the government's problem, with the damage
$d^{\mathrm B}=v'(P)/u'(C)-F_P$ of~\eqref{eq:sp:B:damage} using $\lambda=u'(C)$ from step~(5) and $\tau^Y=0$. The
goods constraint holds by Proposition~\ref{prop:mkt:walras}, the ledger holds by
Definition~\ref{def:mkt:equilibrium}, the state equations are common to the two economies, and the
transversality conditions~\eqref{eq:sp:B:tvc} are the household's and the mine's no-Ponzi conditions.
\end{proof}

As in Version~A, this implements \emph{an} optimum; that it implements \emph{the} optimum needs
uniqueness, which Remark~\ref{rem:sp:B:soc} does not establish.

\begin{remark}[The mass constraint has no market]\label{rem:mkt:B:eta}
Nothing in Definition~\ref{def:mkt:equilibrium} forces $\Phi^IG\le R$. No agent faces the constraint,
and there is no market on which the scarcity of mass for capital formation would show up as a price.
The multiplier $\eta$ is nonetheless implementable with the instruments already listed: it enters as a
uniform reduction of the deposit, $\tau^R=p^W-\eta$, applied to virgin and secondary material alike ---
as it must, since the constraint is on mass and mass is mass whatever its source --- and as an offset
$+\Phi^I\eta$ to the capital rebate, which is capital formation paying for the tonne it competes for.
What is \emph{not} settled is the corner itself: Part~\ref{part:sp} does not characterise the region
where $\eta>0$, so this note does not establish that the equilibrium at the corner is well behaved,
only that the wedge the planner wants there has a market representation. See~\ref{op:sp:corner}
and~\ref{op:mkt:corner}.
\end{remark}


\subsection{Reading the instrument set}\label{sec:mkt:B:reading}

\smalltitle{A deposit--refund scheme at the shadow price of waste.} Away from the corner $\eta=0$, and
the two instruments carrying $p^W$ reduce to a deposit of $\tau^R=p^W$ per tonne of material entering
production and a refund of $s^W=p^W$ per tonne of waste delivered to the collection system. Deposit and
refund are equal, and neither number was chosen: the refund is whatever makes a competitive collection
sector break even, and step~(3) of the proof shows that the planner's waste
condition~\eqref{eq:sp:B:foc:W} \emph{is} that break-even condition. The shadow price of waste is the
break-even collection payment, and the optimal deposit is the same number.

The scheme is source-neutral. Since $R=N+R^R$ the deposit may equivalently be collected from the mine
and from the recycler, and it must be at the same rate: recovered material bears the identical charge
per tonne, because recycling defers waste by one circuit rather than avoiding it
(Section~\ref{sec:sp:B:materials}). A deposit levied on virgin material only, with the proceeds paid
out as a recycling subsidy --- the design most often proposed --- differentiates by source and violates
the arbitrage~\eqref{eq:sp:B:arbitrage} at any rate.

That recycling is nevertheless favoured is priced separately and for physical reasons: the recycler
receives $s^R=d^Wp^P$ because recovery withdraws mass that would otherwise have been emitted at rate
$d^W$, and the mine pays $\tau^{N,S}=\Omega^{N,S}p^W$ because extraction displaces overburden that
recovery does not. These are the two asymmetries between the sources, and they are exactly the two
terms that survive in the arbitrage after $p^W$ cancels.

\smalltitle{The capital instrument is the interest on the deposit.} By~\eqref{eq:mkt:B:tauK} with
$\eta=0$, $\tau^K=-\Phi^I\left(p^W-p^M\right)$: gross capital formation is \emph{subsidised} whenever
$p^W>p^M$. The deposit is charged at entry, on the presumption that the tonne is shed at once. A tonne
embodied in new capital is shed only as the asset depreciates, at a present value of $p^M$ per tonne
by~\eqref{eq:sp:B:costate:M}. The deposit therefore overcharges stored mass by exactly $p^W-p^M$ per
tonne, and $\tau^K$ rebates that on the $\Phi^I$ tonnes carried by each unit of gross formation. Read
this way the rebate is not a subsidy to investment at all but a correction to the timing of a charge
levied elsewhere, and three of its features stop being surprising: its base is gross rather than net
formation, because replacement embodies mass on the same terms as accumulation; it vanishes as
$\delta\to\infty$ and as $r\to0$, by~\eqref{eq:sp:B:q-stationary}, because in either case there is
nothing to postpone; and it flips sign together with $p^W$, since $p^M$ carries the sign of $p^W$
by~\eqref{eq:sp:B:costate:M}.

That last point is worth keeping in view. When waste becomes an asset, $p^W<0$ in the sense
of~\eqref{eq:sp:B:pW-sign}, the deposit turns into a payment for drawing material in, the refund turns
into a charge for handing waste over, and the capital rebate turns into a tax. The three flips are
simultaneous and share one cause; a numerical implementation must leave all three unbounded in sign,
and a solution in which they do not coincide is wrong.

\smalltitle{What is absent.} There is no tax on output, none on consumption, and none on the final good
spent on extraction, exploration or waste treatment. This is the market face of
Proposition~\ref{prop:sp:B:nowedge}, and it deserves to be stated as a result about policy rather than
about a Lagrangian.

\begin{proposition}[No environmental tax on output or consumption]\label{prop:mkt:B:nowedge}
Under Version~B the first best is implemented with $\tau^Y=\tau^C=\tau^N=\tau^D=\tau^W=0$. No
instrument is levied on any use of the final good, notwithstanding that consumption generates
$\Omega^CC$ tonnes of waste and output $\Omega^YY$ tonnes.
\end{proposition}

The reason is the one given after~\eqref{eq:sp:B:FR}: the mass that becomes $\Omega^CC$ was drawn into
the economy as part of $R$ and was charged the deposit at that point. Charging it again when it is
spent prices the same tonne twice. The waste a consumption good becomes is not caused by consuming it;
it is caused by producing the mass, and the mass has already paid.


\subsection{Why the charge is not levied on waste generation}\label{sec:mkt:B:nogate}

The instrument that the coarse intuition recommends is a fee on the waste each activity generates,
levied at the point of disposal --- ``the polluter pays''. Under Version~B it is exactly the wrong
instrument, and the statement is sharp.

\begin{proposition}[A gate fee on waste generation is distortionary]\label{prop:mkt:B:nogate}
Suppose the collection sector charges a fee $g$ per tonne to whoever hands waste over, on the base
$\Omega\mathcal D$, with $s^W$ adjusting through zero profit. Then the equilibrium carries a wedge
$\Omega\omega^jg$ on every unit of the final good spent on activity $j$ and a wedge
$-\Omega\omega^Yg$ on every unit of output produced. Under Version~B the first best requires all such
wedges to be zero, so it is implemented only at $g=0$; any $g\neq0$ is strictly distortionary, whatever
the remaining instruments.
\end{proposition}

\begin{proof}
Spending one unit of the final good on activity $j$ generates $\Omega\omega^j$ tonnes of the fee base
and producing one unit of output generates $\Omega\omega^Y$, so the effective cost of a unit spent is
$1+\Omega\omega^jg$ and the effective value of a unit produced is $1-\Omega\omega^Yg$. By
Proposition~\ref{prop:mkt:B:nowedge} the optimum has all these factors equal to one, and no other
instrument can undo them: the remaining instruments have bases $R$, $N$, $D$, $W$, emissions and $G$,
none of which is proportional to $C$ or to $Y$ except through the allocation itself.
\end{proof}

The wedge $\Omega\omega^jg$ is not merely suboptimal, it is the artefact of~\eqref{eq:sp:B:spurious-C}
made into policy: the fee makes the economy behave as if consuming less reduced the mass drawn in,
which conservation of mass does not permit. A model that posits $W=\Omega Y$ with $\Omega$ fixed
recommends this instrument at rate $g=p^W$; the process-level accounting says the correct rate is zero
and that the same $p^W$ belongs on the deposit instead. Both designs raise the same revenue in a
stationary economy --- $R$ and $W$ then differ only by $\mathcal N$ --- so they are not told apart by
their revenue, only by which margins they distort.

The contrast with Version~A is instructive and is taken up in Section~\ref{sec:mkt:B:compare}: there a
gate fee \emph{is} part of the optimal set, but at the rate $\zeta$, which prices the storage margin
and not the waste.


\subsection{What the government has to know}\label{sec:mkt:B:information}

\begin{proposition}[The optimal instruments do not depend on the composition of waste]
\label{prop:mkt:B:information}
Under Version~B no relative intensity $\omega^j$, and no value of $\Omega$ or $\Phi^Y$, appears in any
instrument of Proposition~\ref{prop:mkt:B:implementation}. The optimal policy is a function of
extraction and exploration costs, the resource rents, the damage function, the recycling and treatment
technologies, the absolute displacement intensities $\Omega^{N,S}$ and $\Omega^{D,S}$, and $\Phi^I$.
\end{proposition}

This follows from Proposition~\ref{prop:sp:B:irrelevance} together with the explicit
solution~\eqref{eq:sp:B:foc:W-solved} for $p^W$, in which no $\omega^j$ and no $\Omega$ appears. It is
the practical counterpart of an uncomfortable result: the same irrelevance means no observation of the
allocation identifies the $\omega^j$ (see~\ref{op:sp:identify}). Here the sign is reversed. The part of
the accounting that is hardest to measure is also the part the government does not need to know, and
the optimal policy is robust to getting it wrong. Under Version~A neither half of that statement
survives.

\smalltitle{The system.} A competitive equilibrium under Version~B consists of the five state
equations~\eqref{eq:sp:B:states-motion}; the household's two
conditions~\eqref{eq:mkt:setup:household-euler}; the final-good firm's
two~\eqref{eq:mkt:setup:foc:goodsfirm}; the recycler's two~\eqref{eq:mkt:setup:foc:recycler}; the
mine's four~\eqref{eq:mkt:setup:foc:mine}--\eqref{eq:mkt:setup:costate:mine}; the collector's
two~\eqref{eq:mkt:setup:foc:varpi} and~\eqref{eq:mkt:setup:zeroprofit}; market
clearing~\eqref{eq:mkt:setup:clearing}; the goods constraint; the ledger, which here solves in one
division as in~\eqref{eq:sp:setup:ledger-solved}; and the government budget. Given instrument paths it
determines the equilibrium under \emph{any} policy, optimal or not --- the object a counterfactual
needs, and one Part~\ref{part:sp} does not contain. The accounting objects $\Omega$, $\Phi^Y$ and
$\Phi^K$ are computed afterwards from~\eqref{eq:sp:B:Omega} and~\eqref{eq:sp:setup:phiK-motion} and do
not feed back, in the market exactly as in the planner's problem.


\subsection{Comparison with Version A}\label{sec:mkt:B:compare}

Everything in this subsection refers to Section~\ref{sec:mkt:A}, and it is the only part of
Section~\ref{sec:mkt:B} that does.

\smalltitle{What is common.} The market structure is identical, by
Proposition~\ref{prop:mkt:common}, and so is the shape of the answer. Under both closures: the
materials block is a source-neutral deposit--refund scheme whose refund is the break-even payment of a
competitive collection sector; the price of material equals the full marginal social cost of a virgin
tonne, $P^R=p^{N,\mathrm B}$ or $p^N$; the price of treated waste is the marginal recycling yield times
the value of a recovered tonne, $P^T=\alpha B^{\mathrm B}$, and $B^{\mathrm B}$ --- the planner's sufficient statistic for the
recycling block --- is a market price; emissions are taxed at $p^P$ with a recovery credit $d^Wp^P$;
the reserve and exploration margins carry no instrument, being internal to the mine; and gross capital
formation is subsidised at the interest on the deposit embodied in durables. None of this depends on
the closure.

\smalltitle{Where they differ.} Table~\ref{tab:mkt:versions} collects the differences.

\begin{table}[H]
\centering
\caption{The optimal instrument set under the two closures.}
\label{tab:mkt:versions}
\setlength{\tabcolsep}{5pt}
\renewcommand{\arraystretch}{1.05}
\begin{tabular}{@{}p{0.28\textwidth}p{0.31\textwidth}p{0.33\textwidth}@{}}
\toprule
 & Version B ($\Phi^I$ exogenous) & Version A ($\Phi^I=\Omega\phi^I$) \\
\midrule
Instruments required & five, plus residual $s^W$ & ten, plus residual $s^W$; six with
the gate fee of Proposition~\ref{prop:mkt:A:gate} \\
Tax on output, on consumption & none & $\Delta\omega^Y$, $\Delta\omega^C$ \\
Tax on extraction, exploration and treatment \emph{expenditure} & none &
$\Delta\omega^N,\Delta\omega^D,\Delta\omega^W$ \\
Gate fee on waste generation & zero; any fee is distortionary
(Proposition~\ref{prop:mkt:B:nogate}) & $\zeta=\chi\left(p^W-p^M\right)$ \\
Deposit $\tau^R$ & $p^W-\eta$ & $p^W-\zeta$ \\
Refund $s^W$ & $p^W$ & $p^W$, or $p^W-\zeta\Omega\mathcal D/W$ with the gate fee \\
Capital instrument $\tau^K$ & $-\Phi^I\left(p^W-p^M-\eta\right)$ &
$-\Phi^I\left(1-\chi\right)\left(p^W-p^M\right)$ \\
Household's $\lambda$ & $=u'(C)$ & $=u'(C)/\left(1+\Delta\omega^C\right)$ \\
Rental of capital $P^K$ & $F_K$ & $\left(1-\Delta\omega^Y\right)F_K$ \\
Depends on $\Omega$ and the $\omega^j$ & no instrument does & every instrument does \\
Signs that flip with $p^W$ & $\tau^R,s^W,\tau^K$ & $\tau^R,s^W,\tau^K$ and the whole
composition block \\
\bottomrule
\end{tabular}
\end{table}

\smalltitle{The one substantive difference.} Strip away the bookkeeping and the two closures disagree
about a single number: the fee that should be charged for generating a tonne of waste. Under Version~B
it is zero, and any positive fee reproduces the artefactual wedge of~\eqref{eq:sp:B:spurious-C}. Under
Version~A it is $\zeta=\chi\left(p^W-p^M\right)$, and a fee is genuinely called for --- not because the
tonne is costly, which the deposit has already handled, but because the spending that generated it
diluted the rate at which the economy stores mass in durables. Neither closure ever justifies the rate
the reduced-form intuition suggests, $g=p^W$. The disagreement is bounded above by
$\chi\left(p^W-p^M\right)$, so it is small when the economy embodies little of its material throughput
in durables and when postponement is worth little, and $\chi$ should be reported alongside $a\varpi$
for exactly this reason (Remark~\ref{rem:sp:coincide}).

\smalltitle{What the choice of closure decides, in the market.} Two things. The first is whether the
optimal policy can be computed without the material-flow composition: under Version~B it can
(Proposition~\ref{prop:mkt:B:information}), under Version~A it cannot, and this is a stronger practical
statement than the identification asymmetry of Section~\ref{sec:sp:B:compare} --- there the $\omega^j$
were merely unidentified, here they are needed to set the tax rates. The second is how many instruments
the policy problem has. Version~B's menu is short enough to be recognisable in existing policy:
a materials deposit--refund, an emission tax with a recycling credit, an extraction charge for
displaced overburden, and an investment subsidy. Version~A adds a waste-generation fee at a rate that
no existing scheme is set at. Neither is adopted as the baseline here, and
Section~\ref{sec:mkt:schemes} classifies the available simplifications under both.
